\begin{ouverture}[courbe]{9}{Matrices et systèmes linéaires}{Graphes et matrices}{Dimensions · Produit · Inverse · Transformations}
\objectif{Lire et calculer avec des matrices de tailles compatibles.}
\objectif{Déterminer un inverse en ordre $2$ ou $3$.}
\objectif{Résoudre un système sous forme matricielle.}
\objectif{Représenter une transformation du plan par une matrice.}
\prerequis{Systèmes linéaires, coordonnées et calcul algébrique.}
\end{ouverture}
\essentiel

\essSection{1}{Un tableau dont les dimensions comptent}
Une matrice de taille $m\times n$ comporte $m$ lignes et $n$ colonnes. Le coefficient $a_{ij}$ est situé à la ligne $i$, colonne $j$.
Une matrice ligne a une seule ligne ; une matrice colonne, une seule colonne. Une matrice carrée d'ordre $n$ a $n$ lignes et $n$ colonnes.
L'addition et la multiplication par un réel se font coefficient par coefficient ; l'addition exige des dimensions identiques.
\begin{definition}[Produit]
Si $A$ est de taille $m\times n$ et $B$ de taille $n\times p$, $AB$ est de taille $m\times p$ et
$(AB)_{ij}=\sum_{k=1}^na_{ik}b_{kj}$.
Chaque coefficient est le produit d'une ligne par une colonne.
\end{definition}
\figurelivre{e09-produit.pdf}{Le coefficient $(1,2)$ de $AB$ se calcule avec la première ligne de $A$ et la deuxième colonne de $B$. La somme des produits croisés vaut ici $1\times1+2\times0=1$.}
\begin{exemple}
Avec $A=\begin{pmatrix}1&2\\3&4\end{pmatrix}$ et $B=\begin{pmatrix}0&1\\1&0\end{pmatrix}$,
$AB=\begin{pmatrix}2&1\\4&3\end{pmatrix}$ mais $BA=\begin{pmatrix}3&4\\1&2\end{pmatrix}$.
En général $AB\ne BA$, même lorsque les deux produits sont définis.
\end{exemple}
\essSection{2}{L'identité et l'inverse}
La matrice identité $I_n$ a des $1$ sur la diagonale et des $0$ ailleurs. Pour une matrice carrée $A$, $AI_n=I_nA=A$ et $A^0=I_n$.
Une matrice $A$ est inversible s'il existe $A^{-1}$ tel que $AA^{-1}=A^{-1}A=I_n$.
\begin{propriete}[Ordre $2$]
Pour $A=\begin{pmatrix}a&b\\c&d\end{pmatrix}$, si $ad-bc\ne0$,
$A^{-1}=\dfrac1{ad-bc}\begin{pmatrix}d&-b\\-c&a\end{pmatrix}$.
Si $ad-bc=0$, $A$ n'est pas inversible. La formule se contrôle en multipliant les matrices.
\end{propriete}
On ne « divise » pas par une matrice ; on multiplie par son inverse, du bon côté, lorsqu'il existe.
\essSection{3}{Systèmes et inverse en ordre $3$}
Le système $2x+y=5$, $x+y=3$ devient
$AX=Y$, avec $A=\begin{pmatrix}2&1\\1&1\end{pmatrix}$, $X=\binom xy$, $Y=\binom53$.
Comme $A^{-1}=\begin{pmatrix}1&-1\\-1&2\end{pmatrix}$, on obtient $X=A^{-1}Y=\binom21$.
En ordre $3$, on peut résoudre successivement les systèmes $AX=e_j$ pour construire les colonnes de l'inverse, ou employer un outil en contrôlant le produit.
\begin{exemple}
$T=\begin{pmatrix}1&1&0\\0&1&1\\0&0&1\end{pmatrix}$ a pour inverse
$\begin{pmatrix}1&-1&1\\0&1&-1\\0&0&1\end{pmatrix}$.
En effet $TX=Y$ donne successivement $x_3=y_3$, $x_2=y_2-y_3$, $x_1=y_1-y_2+y_3$.
\end{exemple}
\essSection{4}{Transformations du plan}
La matrice $R_\theta=\begin{pmatrix}\cos\theta&-\sin\theta\\\sin\theta&\cos\theta\end{pmatrix}$ agit sur la colonne des coordonnées et représente une rotation de centre $O$.
$\begin{pmatrix}1&0\\0&-1\end{pmatrix}$ représente la symétrie par rapport à l'axe horizontal ; $kI_2$ représente l'homothétie de centre $O$ et de rapport $k$.
\figurelivre{e09-transformation.pdf}{La rotation d'angle $\pi/2$ envoie le point $(2,1)$ sur $(-1,2)$. Les vecteurs colonnes et le calcul complexe par multiplication par $\mathrm i$ décrivent la même transformation.}
Une translation non nulle ne se représente pas par une matrice $2\times2$ appliquée seule à la colonne $(x,y)$, car une telle matrice envoie toujours l'origine sur l'origine. Elle s'écrit sous forme affine $X\mapsto AX+C$.

\exercices
\markright{Prendre confiance}
\begingroup\renewcommand{\otBesoinPlace}[1]{}\begin{minipage}{\linewidth}
\rubrique{Prendre confiance}
\begin{exercice}[1]{Dimensions}
$A$ est $2\times3$ et $B$ est $3\times4$.
\begin{questions}
\item Quelle est la taille de $AB$ ?
\item Le produit $BA$ est-il défini ?
\item Peut-on additionner $A$ et $B$ ?
\end{questions}
\end{exercice}
\begin{exercice}[1]{Addition et scalaire}
$A=\begin{pmatrix}1&2\\3&4\end{pmatrix}$.
\begin{questions}
\item Calculer $2A$.
\item Calculer $A+I_2$.
\item Donner $A^0$.
\end{questions}
\end{exercice}
\begin{exercice}[1]{Produit ligne-colonne}
$A=\begin{pmatrix}2&1\\0&3\end{pmatrix}$ et $X=\binom45$.
\begin{questions}
\item Calculer $AX$.
\item Quelle taille a le résultat ?
\item Que représente le premier coefficient ?
\end{questions}
\end{exercice}
\begin{exercice}[1]{Inverser en ordre $2$}
$A=\begin{pmatrix}2&1\\1&1\end{pmatrix}$.
\begin{questions}
\item Calculer $ad-bc$.
\item Donner $A^{-1}$.
\item Vérifier le produit.
\end{questions}
\end{exercice}
\end{minipage}\endgroup
\clearpage
\markright{Appliquer et relier}
\begingroup\renewcommand{\otBesoinPlace}[1]{}\begin{minipage}{\linewidth}
\rubrique{Appliquer et relier}
\begin{exercice}[2]{Non-commutativité}
$A=\begin{pmatrix}1&2\\3&4\end{pmatrix}$ et $B=\begin{pmatrix}0&1\\1&0\end{pmatrix}$.
\begin{questions}
\item Calculer $AB$.
\item Calculer $BA$.
\item Peut-on utiliser $(A+B)^2=A^2+2AB+B^2$ sans condition ?
\end{questions}
\end{exercice}
\begin{exercice}[2]{Résoudre un système}
$2x+y=5$, $x+y=3$.
\begin{questions}
\item Écrire $AX=Y$.
\item Résoudre à l’aide de l’inverse.
\item Vérifier dans les deux équations.
\end{questions}
\end{exercice}
\begin{exercice}[2]{Un ordre $3$}
$T=\begin{pmatrix}1&1&0\\0&1&1\\0&0&1\end{pmatrix}$.
\begin{questions}
\item Résoudre $TX=\begin{pmatrix}3\\5\\2\end{pmatrix}$.
\item Donner $T^{-1}$.
\item Contrôler un coefficient hors diagonale de $TT^{-1}$.
\end{questions}
\end{exercice}
\begin{exercice}[2]{Rotation}
$R=\begin{pmatrix}0&-1\\1&0\end{pmatrix}$.
\begin{questions}
\item Calculer l’image de $(2,1)$.
\item Calculer $R^2$.
\item Interpréter les deux transformations.
\end{questions}
\end{exercice}
\end{minipage}\endgroup
\clearpage
\markright{Chercher et argumenter}
\begingroup\renewcommand{\otBesoinPlace}[1]{}\begin{minipage}{\linewidth}
\rubrique{Chercher et argumenter}
\begin{exercice}[3]{Matrice singulière}
$A=\begin{pmatrix}1&2\\2&4\end{pmatrix}$.
\begin{questions}
\item Pourquoi n’est-elle pas inversible ?
\item Étudier $AX=\binom36$.
\item Étudier $AX=\binom37$.
\end{questions}
\end{exercice}
\begin{exercice}[3]{Ordre d’une composition}
$S=\begin{pmatrix}1&0\\0&-1\end{pmatrix}$ et $R=\begin{pmatrix}0&-1\\1&0\end{pmatrix}$.
\begin{questions}
\item Calculer $RS$ et $SR$.
\item Quelle matrice correspond à « appliquer $S$, puis $R$ » sur une colonne ?
\item Comparer les images de $(1,1)$.
\end{questions}
\end{exercice}
\begin{exercice}[3]{Un inverse de produit}
$A,B$ sont inversibles de même ordre.
\begin{questions}
\item Verifier $(AB)(B^{-1}A^{-1})=I$.
\item Vérifier le produit dans l’autre sens.
\item Conclure et signaler le piège d’ordre.
\end{questions}
\end{exercice}
\begin{exercice}[3]{Translation et matrice}

\begin{questions}
\item Montrer que toute matrice $2\times2$ envoie la colonne nulle sur zéro.
\item Pourquoi une translation non nulle ne peut-elle être représentée ainsi ?
\item Écrire la translation de vecteur $(3,-1)$ sous forme affine.
\end{questions}
\end{exercice}
\end{minipage}\endgroup
